\section{Presentación y alcance}

La práctica estudia riesgos de disponibilidad, confidencialidad e integridad en una LAN Ethernet IEEE 802.3 que transporta IPv4/IPv6 y servicios TCP/IP. El trabajo se realiza exclusivamente en equipos propios o máquinas virtuales, con direcciones privadas, tráfico de prueba y autorización expresa. No se debe generar tráfico hacia Internet, redes institucionales en producción ni direcciones de terceros.

El alcance del informe comprende el análisis de diez casos críticos, la propuesta de controles preventivos, detectivos, correctivos y de recuperación, y el diseño de una solución de programación para el rastreo autorizado de ubicaciones. Los controles se contrastan con las recomendaciones para sistemas de detección y prevención de intrusiones \parencite{nist80094} y para el endurecimiento de servidores \parencite{nist800123}. Cada afirmación técnica debe sostenerse en evidencia verificable: capturas, registros, filtros, comandos, hashes y marcas de tiempo.

\section{Objetivos}

\begin{itemize}[leftmargin=*, topsep=2pt, itemsep=2pt]
	\item Reconocer indicadores de compromiso y comportamientos anómalos en tráfico TCP, UDP, ICMP, DNS, HTTP/HTTPS y mensajería.
	\item Usar Wireshark y registros del sistema para correlacionar eventos, direcciones, puertos, marcas de tiempo y estados de conexión.
	\item Evaluar controles preventivos y detectivos: firewall, segmentación, endurecimiento de servicios, TLS, limitación de tasa e IDS/SIEM.
	\item Proponer el uso responsable de aprendizaje automático para priorizar alertas, sin delegar bloqueos automáticos críticos a un modelo no validado.
	\item Documentar resultados reproducibles mediante capturas, filtros, registros, hashes y recomendaciones verificables.
	\item Diseñar, en la Parte~II, una propuesta técnica y ética de una aplicación, programa o bot que permita compartir y visualizar ubicaciones con consentimiento informado.
\end{itemize}

\section{Resultados de aprendizaje}

\begin{itemize}[leftmargin=*, topsep=2pt, itemsep=2pt]
	\item Diferencia entre síntoma, indicador, hipótesis y evidencia confirmada.
	\item Interpreta capturas de paquetes y registros sin confundir anomalías con prueba definitiva de ataque.
	\item Relaciona cada caso con controles preventivos, detectivos, correctivos y de recuperación.
	\item Entrega un informe técnico con hallazgos, limitaciones y un plan de mejora priorizado.
\end{itemize}

\section{Topología y reglas del laboratorio}

La topología sugerida es un ejemplo; se pueden usar direcciones privadas equivalentes, siempre que la red permanezca aislada y bajo control del equipo.

\begin{table}[H]
	\caption{Elementos de la topología sugerida}
	\label{tab:topologia}
	\centering
	\footnotesize
	\renewcommand{\arraystretch}{1.2}
	\begin{tabularx}{\textwidth}{@{}>{\raggedright\arraybackslash}p{0.25\textwidth} >{\raggedright\arraybackslash}p{0.24\textwidth} Y@{}}
		\toprule
		\textbf{Elemento} & \textbf{Dirección / función} & \textbf{Recomendación} \\
		\midrule
		Switch Ethernet IEEE 802.3 & VLAN de laboratorio & Sin conexión a la red productiva. \\
		Equipo analista & 192.168.50.10/24 & Wireshark; acceso a los registros del laboratorio. \\
		Servidor web/pruebas & 192.168.50.20/24 & HTTP de laboratorio o HTTPS con certificado de prueba. \\
		Servidor DNS de laboratorio & 192.168.50.53/24 & Resolver solo nombres de prueba. \\
		Servidor de mensajería (opcional) & 192.168.50.30/24 & XMPP en entorno aislado. \\
		Gateway de laboratorio & 192.168.50.1/24 & Sin reenvío hacia Internet; filtrado habilitado. \\
		\bottomrule
	\end{tabularx}
\end{table}

\begin{figure}[H]
	\centering
	\fbox{%
		\parbox[c][5.2cm][c]{0.92\textwidth}{%
			\centering
			\textbf{[Insertar aquí el diagrama de la topología del laboratorio]}\\[0.4cm]
			\footnotesize
			Incluya switch, equipo analista, servidores, gateway y las direcciones privadas utilizadas.\\
			Puede elaborarse con draw.io, Diagrams.net, Visio, TikZ o una captura de Packet Tracer.
		}%
	}
	\caption{Topología del laboratorio}
	\label{fig:topologia}
\end{figure}

\noindent\textbf{Reglas obligatorias.}

\begin{itemize}[leftmargin=*, topsep=2pt, itemsep=2pt]
	\item Usar exclusivamente equipos o máquinas virtuales bajo control del grupo y una VLAN o red virtual aislada.
	\item Generar tráfico de prueba benigno o utilizar PCAP didácticos o sintéticos. Ejecutar pruebas de DoS, exfiltración o spoofing únicamente en ambientes aislados y con autorización.
	\item Sincronizar los relojes de los equipos y registrar fecha, hora, IP, puerto, herramienta y versión.
	\item Anonimizar datos personales y credenciales; no capturar tráfico de usuarios ajenos al laboratorio.
	\item Guardar una copia original de cada evidencia y calcular su hash SHA-256; analizar una copia de trabajo.
\end{itemize}

\section{Herramientas sugeridas}

Se sugiere Wireshark; \texttt{netstat} y PowerShell en Windows; el Visor de eventos; Windows Defender Firewall; Zeek y/o Suricata cuando estén disponibles; Wazuh o un SIEM equivalente; y Python con scikit-learn para el análisis de características agregadas \parencite{wireshark,sklearn}. Utilice herramientas y versiones aprobadas para el laboratorio y registre la versión exacta en cada ficha.

\begin{table}[H]
	\caption{Herramientas del laboratorio y registro de versiones}
	\label{tab:herramientas}
	\centering
	\footnotesize
	\renewcommand{\arraystretch}{1.25}
	\begin{tabularx}{\textwidth}{@{}>{\raggedright\arraybackslash}p{0.26\textwidth} Y >{\raggedright\arraybackslash}p{0.15\textwidth}@{}}
		\toprule
		\textbf{Herramienta} & \textbf{Uso en el taller} & \textbf{Versión registrada} \\
		\midrule
		Wireshark & Captura, filtrado y análisis de tráfico. & \\
		\texttt{netstat} y PowerShell & Estado de conexiones, procesos y servicios. & \\
		Visor de eventos & Registros del sistema operativo y de seguridad. & \\
		Windows Defender Firewall & Reglas, registros y validación de filtrado. & \\
		Zeek y/o Suricata & Detección de intrusiones y metadatos de flujos. & \\
		Wazuh o SIEM equivalente & Correlación y priorización de alertas. & \\
		Python y scikit-learn & Análisis de características agregadas y detección de anomalías. & \\
		& & \\
		& & \\
		\bottomrule
	\end{tabularx}
\end{table}

\section{Casos críticos de trabajo}

En cada caso, el equipo debe completar las tareas de observación y análisis, justificar los indicadores, proponer controles y adjuntar evidencia. Las actividades de ataque se sustituyen por análisis de PCAP, telemetría sintética o pruebas benignas autorizadas. La Tabla~\ref{tab:resumen-casos} presenta el panorama de los diez casos.

\begin{table}[H]
	\caption{Resumen de los casos críticos}
	\label{tab:resumen-casos}
	\centering
	\footnotesize
	\renewcommand{\arraystretch}{1.15}
	\begin{tabularx}{\textwidth}{@{}c >{\raggedright\arraybackslash}p{0.27\textwidth} >{\raggedright\arraybackslash}p{0.14\textwidth} Y@{}}
		\toprule
		\textbf{Caso} & \textbf{Tema} & \textbf{Protocolo / servicio} & \textbf{Evidencia principal} \\
		\midrule
		1 & TCP SYN Flood & TCP & Gráfica temporal SYN/ACK y tabla de conexiones. \\
		2 & DNS tunneling & UDP/53, DNS & Consultas filtradas y tabla por host/dominio. \\
		3 & Servicios HTTP, FTP y heredados & HTTP, FTP, otros & Inventario de servicios y matriz de riesgo. \\
		4 & Detección de anomalías con IA/ML & TCP/IP & Notebook reproducible y métricas del modelo. \\
		5 & Resets TCP e indicios de inyección & TCP & Tabla de flujos afectados y extractos de logs. \\
		6 & Reflexión y amplificación UDP & DNS, NTP, SSDP & Gráficas de paquetes y bytes por segundo. \\
		7 & ICMP: escaneo, payloads y túnel & ICMP & Dashboard por IP y tabla de tamaños/frecuencias. \\
		8 & Escaneo de puertos y fingerprinting & TCP/UDP & Lista de exposición y regla de detección. \\
		9 & Priorización de alertas con IA y SIEM & Múltiples fuentes & Pipeline y ejemplo de alerta enriquecida. \\
		10 & Mensajería XMPP y errores de TLS & XMPP, TLS & Checklist TLS y eventos anonimizados. \\
		\bottomrule
	\end{tabularx}
\end{table}
